\section{Discussion}
\label{sec:discussion}

In aggregate, the two paradigms tie on maximum performance. The LLM advantage is narrow and specific: reasoning-heavy retrieval, where bi-encoders score systematically low \citep{su2024bright}.
Pro leads on 5/6 retrieval datasets, showing a consistent advantage across the suite.
Capability profiles show this split: LLMs lead on retrieval while embeddings provide balanced, often superior, coverage across the other four categories (Appendix~\ref{app:figures}, Figure~\ref{fig:radar}).

Viewed through the lens of information retrieval, our results reproduce the classic bi-encoder versus cross-encoder tradeoff at LLM scale~\citep{reimers2019sentencebert, muennighoff2022sgptgptsentenceembeddings}.
One design choice separates the four architectures we evaluate (Figure~\ref{fig:architecture}): how many documents the model is allowed to read jointly with the query.
Embedding pipelines encode queries and documents independently: each document is processed once, offline, and each query requires only vector comparisons.
A cross-encoder reranker reads one document at a time, an LLM listwise reranker reads the top-$k$ shortlist at once, and an LLM in corpus-in-context mode is the extreme case, reading the entire corpus in a single forward pass.
Cost follows that ordering because query-specific computation repeats for each query: reranking scales with the shortlist $k$, and corpus-in-context processing scales with the corpus $N$.
Quality follows the same ordering, which is why reranking is a more economical way to add reasoning than placing the full corpus in context (\S\ref{sec:rerank}).

\begin{figure}[t]
\centering
\includegraphics[width=0.95\textwidth]{architecture_schematic.pdf}
\caption{\textbf{How many documents each architecture reads jointly with the query.}
\emph{Top:} the pipeline. \emph{Bottom:} its attention mask, drawn over the same $N$ documents at the same scale in every panel so the four are directly comparable.
Each row is a token and each column a token it may read.
The red region is what one forward pass reads jointly with the query; it cannot be computed before the query arrives, so cost grows with it.}
\label{fig:architecture}
\end{figure}

Our retrieve-then-rerank experiment (\S\ref{sec:rerank}) shows the same pattern in a production-style setting: a reranker helps most on reasoning-heavy retrieval and little on semantic retrieval.
The same-hardware results attribute much of the quality--cost gap to the underlying inference procedures.

The classification pipelines receive different supervision. Embedding classifiers use kNN over the full labelled training distribution; the LLM receives label names, task descriptions, and up to five examples. Fine-grained tasks make this difference especially visible (\S\ref{sec:fewshot}).
Leading embedding models are also contrastively trained on data overlapping these domains and label spaces, so part of the gap reflects in-domain fit rather than a general representational advantage.
The comparison is therefore between deployment pipelines as practitioners meet them, not intrinsic ceilings: a classification-post-trained LLM would likely narrow the gap, and none was available among the frontier models we evaluate.

The thinking-token tax reveals a mismatch between default reasoning budgets and these tasks.
Reasoning improves the overall Flash score, while lower reasoning budgets preserve retrieval quality for most evaluated models.
Where direct reading comprehension suffices, the model reconsiders straightforward relevance judgements, making the default reasoning budget wasteful for most of these benchmark tasks \citep{lu2025cothurts}.
For practitioners, reducing reasoning effort preserves retrieval quality for most evaluated models while removing the largest component of LLM cost.

\paragraph{Deployment implications.}
These findings motivate a hybrid architecture: embedding models provide high-throughput candidate retrieval, followed by LLM reasoning over a shortlist (Figure~\ref{fig:throughput}).
This division of labour supports retrieve-then-reason pipelines \citep{lewis2020rag, gao2024retrievalaugmentedgenerationlargelanguage}; our reranker results (\S\ref{sec:rerank}) show where they help.
Across the four non-retrieval categories, small-to-medium embedding models closely match the best LLM at a fraction of the cost.
Hardware improvements will reduce the absolute cost of both paradigms. Their different inference procedures preserve an embedding throughput advantage for similarity, classification, and clustering workloads (\S\ref{sec:throughput}).
Accuracy-only leaderboards obscure large cost differences between similarly scored systems. We recommend reporting Pareto frontiers and significance tests alongside accuracy \citep{card2020reporting, dehghani2021benchmarklottery}.

\paragraph{Limitations.}
\label{sec:limitations}
Our ten LLMs from six families are a snapshot of a fast-moving frontier; because \mteblm{} uses the MTEB framework, future models can be evaluated with the same pipeline.
The corpus-in-context protocol places corpora of 82--415 documents entirely in the prompt, which is only feasible at small scale: production corpora require an indexed first stage, so the cost gap we report is a lower bound.
We complement this setting with a bi-encoder, cross-encoder reranker, and LLM listwise reranker comparison on BEIR and BRIGHT (\S\ref{sec:rerank}).
Embeddings also benefit from post-hoc threshold optimisation in pair classification (MTEB convention) with no LLM analogue.
Appendix~\ref{app:extended_limitations} gives the full set.
